\chapter{Integración del intérprete SCPI con el sistema}
\label{cp:integracion-scpi}

El intérprete de comandos SCPI constituye la capa de aplicación del subsistema
de comunicación. Su función es recibir mensajes desde el \textit{middleware},
interpretarlos conforme al estándar IEEE~488.2 y ejecutar las acciones
correspondientes sobre el instrumento. Este capítulo describe cómo el intérprete
—basado en la biblioteca \texttt{libscpi}— se integra con las tres interfaces
que el sistema ofrece: la API de mensajería con el \textit{middleware}, la API
de control del sistema y las APIs de los subsistemas de medición y control.

El capítulo no enumera los comandos SCPI disponibles, cuyo detalle se ha
presentado en el capítulo anterior, sino que expone los mecanismos de conexión
entre el intérprete y el resto del \textit{firmware}.

\section{Integración con la capa de \textit{middleware}}

La comunicación entre el \textit{middleware} y la capa de aplicación se
establece mediante la interfaz abstracta \texttt{app\_msg\_iface\_t}, cuyo
contrato se definió en el capítulo del Sistema de Comunicación. El intérprete
SCPI proporciona una implementación de esta interfaz y la registra en la
\textit{Communication Task} durante la fase de inicialización del sistema.

\subsection{Recepción de bytes crudos}

El \textit{middleware} no realiza delimitación ni análisis de los datos
recibidos. En cada ciclo de la \textit{Communication Task}, los bytes
disponibles en el \textit{buffer} circular de recepción se transfieren
directamente a la aplicación mediante el \textit{callback}
\texttt{on\_rx\_data}. La implementación de este \textit{callback} en el
intérprete SCPI invoca la función \texttt{SCPI\_Input} de \texttt{libscpi}:

\begin{verbatim}
static void on_rx_data(void *ctx, const uint8_t *data, size_t len) {
    (void)ctx;
    SCPI_Input(&scpi_context, (const char *)data, len);
}
\end{verbatim}

\texttt{libscpi} gestiona internamente su propio \textit{buffer} de entrada y
aplica las reglas de terminación configuradas. El intérprete no impone ninguna
estructura de mensajes: la delimitación, el descarte de caracteres de control y
el análisis sintáctico son responsabilidad exclusiva de la biblioteca.

\subsection{Transmisión de respuestas}

Cuando \texttt{libscpi} genera una respuesta (ya sea ante una consulta, un error
o un comando), invoca el \textit{callback} \texttt{write} de su interfaz. La
implementación de este \textit{callback} utiliza la función
\texttt{comm\_app\_send\_response}, que la \textit{Communication Task} expone
como servicio:

\begin{verbatim}
static size_t scpi_write_cb(scpi_t *ctx, const char *data, size_t len) {
    (void)ctx;
    comm_app_send_response((const uint8_t *)data, len);
    return len;
}
\end{verbatim}

La transmisión de la respuesta es asíncrona: los datos se copian al
\textit{buffer} de transmisión de la tarea y la FSM de Transmisión se encarga
del envío hardware. Cuando la transmisión finaliza, el \textit{callback}
\texttt{on\_tx\_done} notifica al intérprete, que puede liberar recursos
asociados a la respuesta.

\subsection{Notificación de errores de comunicación}

Los errores detectados por el \textit{middleware} (desbordamiento del
\textit{buffer} de recepción, fallos de transmisión) se notifican al intérprete
mediante el \textit{callback} \texttt{on\_error}. La implementación actual
registra el evento para su posible consulta mediante el comando
\texttt{SYSTem:ERRor?}, aunque no se toma acción correctiva automática.

\subsection{Secuencia de inicialización}

Durante el arranque del sistema, el módulo \texttt{task\_system} invoca
\texttt{scpi\_engine\_init()} para configurar \texttt{libscpi} con la tabla de
comandos y los \textit{buffers} de entrada. A continuación, obtiene la interfaz
\texttt{app\_msg\_iface\_t} mediante \texttt{scpi\_engine\_get\_iface()} y la
registra en la \textit{Communication Task} con
\texttt{comm\_task\_register\_app\_iface()}. Solo después de este registro el
sistema publica el comando \texttt{COMM\_SYS\_CMD\_START}, asegurando que el
intérprete esté listo para recibir datos antes de que se active la recepción.

\section{Interacción con el sistema}

El intérprete SCPI no actúa de forma aislada: los comandos que afectan al
estado global del instrumento se traducen en solicitudes hacia el módulo de
sistema, y el intérprete es a su vez receptor de notificaciones sobre el estado
del subsistema de comunicación.

\subsection{Comandos que afectan al estado global}

\paragraph{Reset del instrumento (\texttt{*RST}).}
El comando \texttt{*RST} no modifica directamente los parámetros del
instrumento. El intérprete invoca el \textit{callback} estándar
\texttt{SCPI\_CoreRst} para la limpieza de registros IEEE~488.2 y, a
continuación, publica el comando de sistema
\texttt{COMM\_SYS\_CMD\_FULL\_RESET} mediante
\texttt{comm\_sys\_post\_command()}. La tarea de sistema
(\texttt{task\_system}) consume este comando en su ciclo siguiente y ejecuta
\texttt{config\_apply\_reset\_values()}, que restaura los valores por defecto
de todos los canales y apaga las salidas de potencia.

\paragraph{Salvado y carga de configuraciones (\texttt{*SAV}, \texttt{*RCL}).}
Los comandos \texttt{*SAV <slot>} y \texttt{*RCL <slot>}, previstos para una
implementación futura, seguirán el mismo patrón de delegación: el intérprete
publicará un comando de sistema con el número de \textit{slot} como argumento, y
será \texttt{task\_system} quien invoque a \texttt{config\_save\_to\_slot()} o
\texttt{config\_load\_from\_slot()}.

\paragraph{Cambio de interfaz de comunicación.}
El comando \texttt{SYSTem:COMMunicate:IFACe} permite seleccionar el medio de
transporte activo. El intérprete traduce la solicitud a una llamada a
\texttt{comm\_sys\_request\_iface\_change()}, que la FSM de Gestión de la
\textit{Communication Task} evalúa en el ciclo siguiente.

\subsection{Recepción de eventos del sistema}

El intérprete SCPI se registra como suscriptor de las notificaciones del
subsistema de comunicación mediante
\texttt{comm\_sys\_register\_notify\_callback()}. Esto le permite reaccionar
ante cambios de estado sin necesidad de consultas periódicas:

\begin{itemize}
    \item \textbf{Cambio de estado}: ante
          \texttt{COMM\_SYS\_EVENT\_STATE\_CHANGED}, el intérprete puede
          inhabilitar temporalmente la aceptación de comandos si el subsistema
          no se encuentra en estado \texttt{ACTIVE}.
    \item \textbf{Error}: ante \texttt{COMM\_SYS\_EVENT\_ERROR}, el error se
          registra en la cola interna de \texttt{libscpi} para su consulta
          mediante \texttt{SYSTem:ERRor?}.
    \item \textbf{Conexión / desconexión}: permiten reiniciar el estado del
          \textit{parser} cuando se establece o se pierde la conexión con el
          \textit{host}.
\end{itemize}

Este mecanismo convierte al intérprete en un suscriptor del sistema, aumentando
la robustez sin introducir acoplamiento adicional.

\section{Interacción con los subsistemas de control}

El instrumento organiza sus recursos de medición y control en \textbf{canales}.
Cada canal agrupa tres componentes —un sensor de temperatura, un controlador
PID y una etapa de salida— bajo una estructura contenedora denominada
\texttt{controller\_t}. Esta estructura se obtiene mediante la función
\texttt{controller\_get\_instance(channel)}, que retorna la instancia
correspondiente al número de canal.

Los \textit{callbacks} de los comandos SCPI no acceden directamente al
hardware, sino que invocan las funciones definidas en las APIs de cada
componente. Esto desacopla el intérprete de la implementación concreta y
permite la evolución independiente de cada subsistema.

\subsection{API de sensor}

Cada sensor de temperatura expone sus servicios a través de la interfaz
\texttt{sensor\_api\_t}.

\begin{table}[htbp]
    \centering
    \caption{Funciones de la API de sensor.}
    \label{tab:sensor-api}
    \begin{tabular}{ll}
        \toprule
        \textbf{Función} & \textbf{Descripción} \\
        \midrule
        \texttt{read\_temperature(sensor)} & Retorna la temperatura actual en °C. \\
        \texttt{set\_type(sensor, type)}   & Configura el tipo de sensor (RTD, TC, etc.). \\
        \texttt{get\_type(sensor)}         & Retorna el tipo de sensor configurado. \\
        \texttt{set\_protection\_min(sensor, min)} & Establece el límite inferior de protección. \\
        \texttt{get\_protection\_min(sensor)}      & Retorna el límite inferior de protección. \\
        \texttt{set\_protection\_max(sensor, max)} & Establece el límite superior de protección. \\
        \texttt{get\_protection\_max(sensor)}      & Retorna el límite superior de protección. \\
        \texttt{enable\_protection(sensor, enable)}& Activa o desactiva la protección. \\
        \texttt{is\_protection\_enabled(sensor)}   & Indica si la protección está habilitada. \\
        \texttt{is\_protection\_tripped(sensor)}   & Indica si la protección se ha disparado. \\
        \texttt{clear\_protection(sensor)}         & Reanuda la protección tras un disparo. \\
        \bottomrule
    \end{tabular}
\end{table}

La protección del sensor opera de forma independiente al control PID. Si la
temperatura medida supera el límite máximo o desciende por debajo del mínimo,
la salida de potencia correspondiente se desactiva automáticamente.

\subsection{API de controlador PID}

El controlador PID proporciona las funciones de sintonía y ejecución del lazo
de control. La API no incluye los límites de salida, que son responsabilidad de
la etapa de potencia.

\begin{table}[htbp]
    \centering
    \caption{Funciones de la API de PID.}
    \label{tab:pid-api}
    \begin{tabular}{ll}
        \toprule
        \textbf{Función} & \textbf{Descripción} \\
        \midrule
        \texttt{set\_setpoint(pid, sp)} & Establece el punto de consigna. \\
        \texttt{get\_setpoint(pid)}     & Retorna el punto de consigna. \\
        \texttt{set\_kp(pid, kp)}       & Establece la ganancia proporcional. \\
        \texttt{get\_kp(pid)}           & Retorna la ganancia proporcional. \\
        \texttt{set\_ki(pid, ki)}       & Establece la ganancia integral. \\
        \texttt{get\_ki(pid)}           & Retorna la ganancia integral. \\
        \texttt{set\_kd(pid, kd)}       & Establece la ganancia derivativa. \\
        \texttt{get\_kd(pid)}           & Retorna la ganancia derivativa. \\
        \texttt{compute(pid, input)}    & Ejecuta una iteración del PID. \\
        \texttt{reset(pid)}             & Reinicia el estado del PID. \\
        \bottomrule
    \end{tabular}
\end{table}

\subsection{API de salida}

La etapa de salida gestiona el ciclo de trabajo (\textit{duty cycle}) y los
límites que restringen su excursión. También controla la activación y
desactivación de la potencia.

\begin{table}[htbp]
    \centering
    \caption{Funciones de la API de salida.}
    \label{tab:output-api}
    \begin{tabular}{ll}
        \toprule
        \textbf{Función} & \textbf{Descripción} \\
        \midrule
        \texttt{set\_state(out, on)}       & Activa o desactiva la salida. \\
        \texttt{get\_state(out)}           & Retorna el estado de la salida. \\
        \texttt{set\_limit\_min(out, min)} & Establece el límite inferior de excursión. \\
        \texttt{get\_limit\_min(out)}      & Retorna el límite inferior. \\
        \texttt{set\_limit\_max(out, max)} & Establece el límite superior de excursión. \\
        \texttt{get\_limit\_max(out)}      & Retorna el límite superior. \\
        \texttt{get\_duty\_cycle(out)}     & Retorna el ciclo de trabajo actual. \\
        \texttt{is\_limited(out)}          & Indica si la salida está saturada en un límite. \\
        \bottomrule
    \end{tabular}
\end{table}

\subsection{Mapeo de comandos SCPI a las APIs de control}

La tabla \ref{tab:scpi-to-api} resume la correspondencia entre los comandos
SCPI definidos para el instrumento y las funciones de las APIs que se invocan
para atenderlos. En todos los casos, el número de canal se obtiene mediante
\texttt{SCPI\_CommandNumbers(context, \&channel, 1, 1)} y la instancia del
controlador se recupera con \texttt{controller\_get\_instance(channel)}.

\begin{table}[htbp]
    \centering
    \caption{Correspondencia entre comandos SCPI y APIs de control.}
    \label{tab:scpi-to-api}
    \begin{tabular}{ll}
        \toprule
        \textbf{Comando SCPI} & \textbf{Función invocada} \\
        \midrule
        \texttt{MEASure:TEMPerature\#?} & \texttt{sensor\_api\_t::read\_temperature} \\
        \texttt{SENSe\#:TEMPerature:TYPE} & \texttt{sensor\_api\_t::set\_type} \\
        \texttt{SENSe\#:TEMPerature:TYPE?} & \texttt{sensor\_api\_t::get\_type} \\
        \texttt{SENSe\#:TEMPerature:PROTection:MIN} & \texttt{sensor\_api\_t::set\_protection\_min} \\
        \texttt{SENSe\#:TEMPerature:PROTection:MIN?} & \texttt{sensor\_api\_t::get\_protection\_min} \\
        \texttt{SENSe\#:TEMPerature:PROTection:MAX} & \texttt{sensor\_api\_t::set\_protection\_max} \\
        \texttt{SENSe\#:TEMPerature:PROTection:MAX?} & \texttt{sensor\_api\_t::get\_protection\_max} \\
        \texttt{SENSe\#:TEMPerature:PROTection:STATe} & \texttt{sensor\_api\_t::enable\_protection} \\
        \texttt{SENSe\#:TEMPerature:PROTection:STATe?} & \texttt{sensor\_api\_t::is\_protection\_enabled} \\
        \texttt{SENSe\#:TEMPerature:PROTection:TRIPped?} & \texttt{sensor\_api\_t::is\_protection\_tripped} \\
        \texttt{SOURce\#:CONTrol:SETPoint} & \texttt{pid\_api\_t::set\_setpoint} \\
        \texttt{SOURce\#:CONTrol:SETPoint?} & \texttt{pid\_api\_t::get\_setpoint} \\
        \texttt{SOURce\#:CONTrol:PID:KP} & \texttt{pid\_api\_t::set\_kp} \\
        \texttt{SOURce\#:CONTrol:PID:KP?} & \texttt{pid\_api\_t::get\_kp} \\
        \texttt{SOURce\#:CONTrol:PID:KI} & \texttt{pid\_api\_t::set\_ki} \\
        \texttt{SOURce\#:CONTrol:PID:KI?} & \texttt{pid\_api\_t::get\_ki} \\
        \texttt{SOURce\#:CONTrol:PID:KD} & \texttt{pid\_api\_t::set\_kd} \\
        \texttt{SOURce\#:CONTrol:PID:KD?} & \texttt{pid\_api\_t::get\_kd} \\
        \texttt{SOURce\#:CONTrol:OUTPut?} & \texttt{output\_api\_t::get\_duty\_cycle} \\
        \texttt{SOURce\#:CONTrol:OUTPut:LIMit:MAX} & \texttt{output\_api\_t::set\_limit\_max} \\
        \texttt{SOURce\#:CONTrol:OUTPut:LIMit:MAX?} & \texttt{output\_api\_t::get\_limit\_max} \\
        \texttt{SOURce\#:CONTrol:OUTPut:LIMit:MIN} & \texttt{output\_api\_t::set\_limit\_min} \\
        \texttt{SOURce\#:CONTrol:OUTPut:LIMit:MIN?} & \texttt{output\_api\_t::get\_limit\_min} \\
        \texttt{SOURce\#:CONTrol:OUTPut:LIMit:STATe?} & \texttt{output\_api\_t::is\_limited} \\
        \texttt{OUTPut\#:STATe} & \texttt{output\_api\_t::set\_state} \\
        \texttt{OUTPut\#:STATe?} & \texttt{output\_api\_t::get\_state} \\
        \bottomrule
    \end{tabular}
\end{table}

\section{Consideraciones sobre la implementación actual}

En el estado actual del \textit{firmware} (junio de 2026), las APIs de sensor,
PID y salida descritas en este capítulo constituyen un modelo de referencia
hacia el cual evolucionará el código. La implementación presente utiliza una
única instancia de controlador PID accesible mediante
\texttt{get\_temp\_pid\_instance()} y carece de la estructura contenedora
\texttt{controller\_t} y de las APIs de sensor y salida como entidades
independientes. Los \textit{callbacks} del intérprete SCPI invocan directamente
las funciones del módulo \texttt{pid\_controller} y del \texttt{ADC} para la
lectura de temperatura.

Esta simplificación es adecuada para la fase actual de desarrollo, que prioriza
la funcionalidad sobre la abstracción. La migración hacia el modelo descrito en
este capítulo se realizará de manera progresiva, sin modificar la interfaz
pública del intérprete SCPI ni la de comunicación con el \textit{middleware}.